\subsection{Pcontrol: Scheduling a Recursive Object Computation}
\label{sec:pcontrol}

The programs in \filepath{examples/pcontrol/} move reflection from an
individual object's behavior to a group's execution policy. A recursive
algorithm can create many objects before earlier work has finished.
The example asks whether choosing small pending computations first can
reduce this growth. Here, \emph{parallelism control} means controlling the
progress of concurrent work, not creating operating-system threads.
Read the application first, then follow the machinery that schedules it.

\subsubsection{A Checked Session and the Loading Order}
\label{sec:pcontrol-load}

Start a fresh interpreter with its working directory set to a scratch
copy of \filepath{examples/pcontrol/}. Relative names in
\filepath{start.abclr2} refer to that directory. Loading creates translations
and compiling creates compiled files; the experiments also overwrite
their own log files. The automated reading check creates this scratch
copy for you:

\begin{lstlisting}
make -C tutorial pcontrol-check
\end{lstlisting}

For an interactive session in such a copy, use:

\begin{codeblock}
\begin{lstlisting}
(inc "start" :use-source t)
(inc "fib" :use-source t)
[test-pc <== [:start 4]]
; Prints Fibonacci(4) = (:ANS 3); replies T.
[test-np <== [:start 4]]
; Prints Fibonacci(4) = (:ANS 3); replies T.
\end{lstlisting}
\end{codeblock}

The printed \code{(:ANS 3)} is the Fibonacci answer message, whereas
\code{T} is the test object's final reply. These are different protocol
results. Use the argument shapes in \filepath{README.abclr2} and the
actual handlers: some commands in the commented block at the end of
\filepath{start.abclr2} have obsolete extra arguments.
Keep inputs small; the original README warns that these experiments
are expensive.

\filepath{start.abclr2} compiles both applications but loads only the
four machinery files. Their order supplies the dependencies:

\begin{enumerate}[nosep]
  \item \filepath{para-monitor.abclr2} defines the monitor.
  \item \filepath{pqueue-eval.abclr2} defines the queue and evaluator factories.
  \item \filepath{priori-meta.abclr2} defines application metaobjects.
  \item \filepath{pcontrol-group.abclr2} connects them into a group.
\end{enumerate}

The subsequent \code{inc "fib"} or \code{inc "qsort"} chooses the
application. Both application files define \code{test-pc} and
\code{test-np}; loading the second replaces those bindings. Use a fresh
session when changing applications.

\subsubsection{Fibonacci: Create Children, Then Collect Replies}
\label{sec:pcontrol-fib}

The convention is \code{F(1) = F(2) = 1}. The source does not provide a base
case for zero or negative inputs. A factory \code{fib-gen-gen} creates
a \code{fib-gen} in the \code{pcontrol} group. That generator creates
one \code{fib} object per recursive subproblem, with its problem size
as its priority. The child objects inherit the creation context's group.

\wholelisting[firstline=18,lastline=33]{examples/pcontrol/fib.abclr2}

The literal patterns \code{1} and \code{2} precede the general
variable pattern \code{n}. A base case sends \code{[:ans 1]} to its
request's reply destination, then reports logical completion to the
monitor. A larger problem saves its caller's destination in state
\code{reply}. Each nested \code{fib-gen <== [:new ...]} waits for
the newly created child object, not for its Fibonacci result.
The following \code{<=} starts that child asynchronously with
\code{@ Me}, so the child sends its answer back to the parent.

The parent does not block while collecting answers. Its first
\code{[:ans n]} stores the value in \code{sub-node-value}; its
second adds both values and forwards \code{[:ans sum]} to the
saved destination. Because all valid results are positive, \code{nil}
can safely mean no first answer yet. This is a one-computation-per-object
protocol: the state is not reset for a second request.

Trace the five application nodes for input 4:

\begin{enumerate}
  \item The root for 4 saves its caller and creates children for 3 and 2.
  \item The node for 3 saves the root's identity and creates children
        for 2 and 1.
  \item Each base node replies \code{[:ans 1]} to its parent.
  \item The node for 3 combines two answers into \code{[:ans 2]}.
  \item The root combines 2 and 1 into \code{[:ans 3]}.
\end{enumerate}

This is a dependency trace, not a guaranteed temporal ordering.
At each join either child may reply first. Addition makes the
aggregation independent of that order. The test waits on the root,
prints its answer, requests that monitoring stop, and replies \code{T}.
Its \code{<==} is therefore a completion boundary for the computation,
not just for dispatching the first child.

\subsubsection{Where the Scheduling Policy Actually Lives}
\label{sec:pcontrol-policy}

The group definition is short enough to read in full:

\wholelisting[firstline=9,lastline=16]{examples/pcontrol/pcontrol-group.abclr2}

The group uses \code{priori-meta-gen} for each application object's
metaobject. Its evaluator is a \code{multi-eval} created with four
primary evaluators. The group forwards an object-creation request,
including \code{:priority priority}, to its metaobject generator.
Four evaluators here are four runtime objects; this declaration does
not establish execution on four processors.

There are two queues to distinguish. Each application metaobject in
\filepath{priori-meta.abclr2} has its own ordinary FIFO message queue.
On \code{:begin}, it selects a script and submits a \code{[:do ...]}
request to the shared evaluator, adding the object's priority after
the usual context fields. A continuation sends \code{:end} back to
that metaobject when script evaluation completes. That signal advances
the object's queue; it is separate from the script's application reply.

The shared \code{multi-eval} has a metaobject made by \code{PQ-meta-gen}.
That metaobject queues evaluation requests as
\code{[Priority [Message Reply Sender]]}.
An explicitly supplied priority is associated with the requesting
object's \code{Id}; requests without an explicit priority use the
recorded association, or zero if none exists.
The priority is thus scheduling information in the evaluation context,
not an extra field in \code{[:ans value]}.

Both tests keep this machinery and the \code{pcontrol} group.
They replace the shared evaluator metaobject's queue:

\wholelisting[firstline=44,lastline=45]{examples/pcontrol/fib.abclr2}

The \code{pc} variant installs a queue sorted by \code{<} on each
entry's \code{first} element. The \code{np} variant instead uses
\code{[queue-gen <== :new]}, the ordinary FIFO queue.
In the nested expression, \code{[pcontrol <== :evaluator]} obtains
the evaluator object; \code{meta} obtains its metaobject; and
\code{[:set-queue ...]} changes that metaobject's pending-work policy.
This does not change each Fibonacci object's inbox into a priority queue.

\wholelisting[firstline=12,lastline=18]{examples/pcontrol/pqueue-eval.abclr2}

\code{cons} inserts an entry into the list and \code{sort} reorders
that list. \code{:key \#'first} tells sorting to compare the priority
field. Smaller problem sizes are selected before larger ones among
entries pending in this queue. Do not infer a stable order for equal
priorities: the source uses \code{sort}, not a stable tie-breaking
policy. Nor is this express interruption: it selects pending requests
rather than preempting an already running script.

Finally, read \code{multi-eval}'s two handlers near the end of
\filepath{pqueue-eval.abclr2}. One recognizes an object-definition
representation carrying \code{:priority} and forwards creation to
the chosen group manager, then arranges the resulting name binding
and reply. The catch-all handler forwards other evaluation requests
to the next primary evaluator in a cyclic list, preserving their
completion destination \code{C}. This explains how the example
changes creation and scheduling without rewriting the Fibonacci
application's reply protocol.

\subsubsection{Quicksort: A Message-driven State Machine}
\label{sec:pcontrol-qsort}

\filepath{qsort.abclr2} uses the same group but a different application.
Its opening comment explains why scanning is implemented through
self-messages rather than a \code{do} loop. Read those self-messages
as transitions between phases, not as recursive host function calls.

Each sorting object captures \code{array}, \code{left}, and
\code{right} from its generator's lexical scope. Its priority is
\code{(- right left)}. For a nontrivial range, \code{:start}
saves its caller's destination, chooses the middle element as
\code{pivot}, initializes indices \code{i} and \code{j}, and sends
\code{:scan-left} to itself.

\begin{enumerate}
  \item \code{:scan-left} advances \code{i} until its element is
        not smaller than the pivot, then requests \code{:scan-right}.
  \item \code{:scan-right} moves \code{j} left until its element
        is not larger than the pivot, then requests \code{:test-swap}.
  \item \code{:test-swap} exchanges elements if \code{i <= j},
        advances both indices, and either resumes scanning or creates
        two child sorting objects for the remaining ranges.
  \item Each child sends \code{:fin} to its parent. The parent's
        \code{wait} flag records the first completion; the second
        forwards \code{:fin} to the parent's saved caller.
\end{enumerate}

\code{aref} reads an array element; \code{setf} changes that element
in place. The descendants share this array but work on partitioned
ranges. The first completion means only that one side is sorted;
the second establishes completion of both. A range with
\code{left >= right}, including an empty range, immediately replies
\code{:fin}.

\wholelisting[firstline=54,lastline=62]{examples/pcontrol/qsort.abclr2}

For \code{(3 1 3 2)}, the root covers indices 0 through 3 and
chooses the element at index 1, namely 1, as its pivot.
The left scan stops at index 0. The right scan moves from 3 to 1.
The swap produces \code{(1 3 3 2)}, with \code{i = 1} and
\code{j = 0}. Scanning ends, so the children cover the already
trivial range 0--0 and the remaining range 1--3. Their completions
join at the root; the final array is \code{(1 2 3 3)}.

In a fresh session with the same setup, a deterministic test avoids
the random initialization object:

\begin{codeblock}
\begin{lstlisting}
(inc "start" :use-source t)
(inc "qsort" :use-source t)
[input-array := (make-array 4 :initial-contents '(3 1 3 2))]
(coerce [test-pc <== [:start input-array 4]] 'list)
; => (1 2 3 3)
(coerce [test-np <== [:start input-array 4]] 'list)
; => (1 2 3 3)
(coerce input-array 'list)
; => (3 1 3 2)
\end{lstlisting}
\end{codeblock}

\code{coerce ... 'list} makes an array's contents convenient to inspect.
Both test handlers first copy their input with
\code{make-array :initial-contents}; they return the sorted copy,
not the caller's modified array. The original random entry point
\code{[init <= N]} instead stores an array in the outer binding
\code{a}; follow its \code{sample-gen} self-message loop to see
how it fills each index.

\subsubsection{Reading the Monitor Without Overinterpreting It}
\label{sec:pcontrol-monitor}

\filepath{para-monitor.abclr2} records two columns: a logical
\code{clock} and a counter. The evaluator sends \code{:tick}
when handling an evaluation request; this clock is not elapsed
wall-clock time. Application metaobjects announce creation and
transitions between dormant and active dispatch. The application
explicitly sends \code{:object-dead} at logical completion.
Nothing in that message destroys a Lisp object or measures garbage
collection.

The active count tracks announced dispatch activity, not the number
of physical CPUs busy. A node waiting for child answers can be
dormant yet still represent unfinished application work.
The object count combines creation announcements with explicit
logical-death announcements. For example, generator objects are
also created through the machinery. It is not simply a count of
the recursive leaves, allocated bytes, or all reachable host objects.
The \code{:record} messages are asynchronous observations of these
counters.

For Fibonacci input 4, the actual source names the files
\code{fib-active-pc-4.log}, \code{fib-object-pc-4.log}, and
the corresponding \code{np} files; quicksort uses a \code{qsort-}
prefix. Repeating the same run supersedes those files.
The directory README's generic names omit these application prefixes.
Compare traces as experiment output, not as proof that one policy
always reduces memory or is faster.

\paragraph{Check your understanding.}
Does \code{test-np} remove priority-bearing metaobjects?
\emph{Answer:} No. It replaces the shared evaluator's work queue with
a FIFO queue while retaining the group and application metaobjects.
Why can a parent accept either child's answer first?
\emph{Answer:} Its state records the first answer and combines it
with the second; Fibonacci uses addition, while quicksort counts
two completion signals for disjoint subranges.

\paragraph{Reproducing the reading checks.}
\code{make -C tutorial pcontrol-check} runs twelve checks in fresh
sessions using copies of the distributed sources. They check priority-queue
ordering for distinct priorities, both Fibonacci policies at 4 and 6,
both quicksort policies on a duplicate-containing input, and preservation
of the caller's array. Translations, monitor files, inputs, and
transcripts stay in \filepath{build/pcontrol-tests/}.
These results establish the small examples' values and completion,
not general performance, queue fairness, or faithful historical
parallel execution.
